\summarychapter{39}{The Destiny of a Nation}{A Christian America}

\begin{inbrief}
The dream of a Christian America, a ``righteous empire,'' is a dominant theme of nineteenth-century
American Protestantism. The chapter traces its rise on the western frontier through voluntary
societies and revivals, its great crisis over slavery and the Civil War, and its decline under three
cultural shocks: Darwin, the industrial city with its immigrants, and the higher criticism of the
Bible.
\question{How did Evangelicals try to shape the destiny of the American nation, and why did the vision
fade?}
\end{inbrief}

``Shall a nation be born at once?'' asked Isaiah, and in 1835 Lyman Beecher, the well-known
Presbyterian and Congregational minister, preached on the text a sermon called \emph{A Plea for the
West}. A vast empire was opening in the wilderness and a whole culture was at stake; Christians
should shape ``the religious and political destiny of the nation.'' His means were preaching,
distributing Bibles, planting churches, founding schools and reforming morals, for, Puritan that he
was, he knew that a free society needs just laws and that in a democracy just laws need popular
support informed by Christianity. He spoke for a host of Baptists, Methodists, Congregationalists,
Presbyterians and Episcopalians.

\ssection{The Challenge of the American West}
Like Innocent III's vision of medieval Europe and Calvin's design for Geneva, the evangelical dream
was an expression of Christian concern for society, and the West gave it the greatest opportunity
ever to press the gospel on a whole nation. The vision has faded, Shelley writes, but the impulse
remains, since Christ calls his people light and salt. The great fact of the century was the moving
frontier. Early visitors praised the land beyond the Alleghenies, with its streams and meadows of
wild rye, bluegrass and clover. After the Revolution the continent seemed to tilt toward the Pacific:
nine states were added between 1792 and 1821, and by mid-century half the population lived west of
the Appalachians. The frontiersman, driving out the Indians and subduing the wilderness, gained a
name for lawless living; an English visitor described backwoodsmen fighting over trifles with hands,
teeth, knees and feet and gouging out eyes. At the nation's birth the churches seemed ill prepared.
Only 5 or 10 percent of Americans were church members. Yet the crude and godless West was tamed, and
evangelical Christianity did more than any other force to tame it.

It had two instruments. The first was the voluntary society. The Bill of Rights, by guaranteeing
religious liberty, sanctioned the denominational idea of the church and ruled out direct church
influence on government. Carey and the English Evangelicals had devised the voluntary society for
missions and for the fight against the slave trade; Americans took it over to influence public
opinion, support missionary and educational work and spread reform, uniting Christians of different
denominations for a common concern such as temperance or the Sabbath. So arose the American Bible
Society, the American Colonization Society, the American Sunday School Union, the American Education
Society and many more. By 1830 Beecher credited the church's ``grand influence'' to ``the voluntary
association of Christians.''

The second was the revival, to which the churches could look, Beecher said, for members, pastors and
power over public opinion. In 1790 the task was double, to regain the East and to win the West. In
the East, especially in colleges, new spiritual life appeared before 1800, the Second Great
Awakening, which supplied leaders for the western crusade. The frontier revival, centered in
Kentucky and Tennessee, was as rugged, wild and boisterous as the people. Timothy Flint described the
western preacher, made pensive and romantic by months of travel through dark forests, whose first aim
was to stir the feelings with ``strong, earnest and unlettered eloquence.'' Among a largely
illiterate people the frontier preacher, white or black, was the chief channel of Christianity.

\ssection{The Inflammable James McGready}
One of the most pungent was James McGready, a tall, angular Presbyterian of Scotch-Irish stock from
Pennsylvania, with keen black eyes and an uncompromising manner. He first preached in North Carolina,
where his stress on God's wrath lit a revival; when he left abruptly for Kentucky in 1798, it was
said that the suggestion had come in a letter written in blood. In Logan County, a ``Rogues'
Harbor'' of horse thieves and murderers, he served three small congregations at Red River, Gasper
River and Muddy River. He painted heaven so vividly that hardened hearers longed for it, and hell so
that they trembled before the lake of fire. In July 1800 came his Pentecost. After a revival at Red
River he sent out notice of the next sacramental service at Gasper River, and pioneers came by wagon,
on horseback and on foot from as far as a hundred miles, with tents and provisions, ready to stay.
This was the first camp meeting: a religious service of several days held outdoors for people who
had come a distance and camped on the spot.

Revival preacher and camp meeting have lasted in America for nearly two centuries. The heat cooled,
as it had to, and under Charles Finney, D.~L.~Moody and Billy Graham the camp meeting moved indoors
to rural chapels and city auditoriums. Not everyone approved. Many Lutherans and Presbyterians thought
revivals slighted sound doctrine, and Catholics and Episcopalians thought them emotional outbursts,
not worship. The revivalists and missionaries went on preaching, planting churches and founding
colleges, and by the 1830s Tocqueville found no country in the world where Christianity kept ``a
greater influence over the souls of men.''

But a cancer was spreading through Christian America. How could a democracy infused with Christian
principles go on enslaving millions? American slavery began on August 20, 1619, when twenty Africans
were landed from a Dutch ship at Jamestown; by 1830 there were about two million slaves. With each
new state and each move onto new land the question burned deeper: slave or free? If liberty and union
could not both be kept, which came first? The struggle was so basic and so religious that all sides
turned to the Bible.

\ssection{Christianity among the Slaves}
The slave turned to the Bible out of necessity, because his master had stripped him of everything
else, including his African gods. Uprooting had shattered the Africans' culture, their social
organization and their sense of a place in the world. Some killed themselves on the ``middle
passage'' across the Atlantic and others tried to escape, but most submitted and sought a meaning for
life in the white man's world. Christianity gave them a new center. Some owners at first opposed
teaching slaves the Scriptures for fear they would find there ideas of equality and rebel. The
opposition faded as masters decided that the New Testament supports slavery and that the best slaves
were those who knew the Bible best. Since every other form of organized life was forbidden, the black
preacher became the central figure of the ``invisible institution,'' the slave church. He was
``called'' by an experience showing that God had chosen him, which gave him great influence, and he
learned to dramatize the Bible in terms of the slaves' own life. As Dean Willard Sperry explained,
the chief event was the Exodus and the chief character Jesus. Bondage in Egypt, the Red Sea, the
drowning of Pharaoh's army, the desert and the crossing of Jordan into the Promised Land were the
constant themes of sermons and spirituals. Because there was little hope of deliverance on earth,
emancipation was linked with death, when Jesus would strike off the shackles.

Most Christians did not try to defend slavery before the 1830s, and in the first three decades of
the century the antislavery movement was stronger in the South than in the North. Then it faded
there and a Southern defense arose, partly from the Bible. Richard Furman, South Carolina's leading
Baptist, argued that Israel was told to buy bondmen who were to be ``bondmen for ever,'' and that the
apostles never demanded emancipation but only that masters give slaves what is ``just and equal.''
In the years of the Cotton Kingdom the tie between evangelical religion and race became part of the
Southern way of life. The South grew isolated from the North and from a Western world that condemned
slavery, and isolation made it defensive. Harriet Martineau, visiting in the 1830s, said a ``magic
ring'' seemed drawn around those who lived amid slavery, so that few saw anything beyond it.

In the North evangelical feeling against slavery flowed largely from the revival preaching of
Charles G.~Finney, especially around Oberlin College, of which he was president. His disciple
Theodore Weld led the crusade with \emph{The Bible against Slavery} (1837) and \emph{Slavery as It
Is}. Beecher's daughter, Harriet Beecher Stowe, lived with the latter book ``day and night till its
facts crystallized into \emph{Uncle Tom's Cabin}.'' Her novel reminds Christians that prophecy joins
the day of vengeance with the year of the redeemed. In the common millennialist reading of
Revelation the fall of Babylon would be marked by the convulsion of nations, and America, with its
part in God's plan, must be purged of its guilt. The guilt was national, not Southern only, and she
struck at the national conscience hoping that a cleansing of the soul might avert a scourging of the
body.

All sides thus used the same symbols---one Bible, one heaven and hell, one Christ---for opposite
causes. No one felt this more than Abraham Lincoln. Formed by evangelical culture and by the language
of the Bible, he never joined a church or held a particular creed, and he learned that no one can say
precisely what God wills for the nation. ``In great contests each party claims to act in accordance
with the will of God. Both may be, and one must be wrong.'' In his second inaugural address he said
that both sides read the same Bible and pray to the same God; the prayers of both could not be
answered, for ``the Almighty has his own purposes.'' The war was fought and the nation endured. The
vision of a Christian America survived as well, but much weakened. Black churches arose in great
numbers, the chief institutional expression of freedom and a standing reminder of the vision's blind
spot.

\ssection{Cultural Shocks for Evangelical America}
Hardly was the war over when three shocks called traditional evangelical beliefs into question. The
first was Charles Darwin's \emph{Origin of Species} (1859), perhaps the century's most important
book. It argued that species evolve by natural selection: those that adapt to their surroundings
survive and the weaker die off. Would such a theory leave room for a God who creates and sustains,
and how was the biblical story of creation to be read? The second was industrialization and the rush
to the cities. Towns became cities overnight, filled from the countryside and from Germany, Norway,
Italy and elsewhere, and most immigrants brought religious views foreign to the way Protestant
Americans had seen their country and their Bible. The third and most direct was higher criticism. As
professors returned with advanced degrees from European universities, critical views came to dominate
American higher education and then many denominations. Churches were jolted to hear that Moses did
not write the first five books and that Jesus was a somewhat deluded visionary.

Together these belonged to the general shift of Western culture from Christian to secular thought and
conduct, and Christians disagreed about the response. In the evangelical denominations two parties
formed, one welcoming the changes as blessings from God, the other resisting them as threats to the
biblical message. Dwight L.~Moody, the leading urban revivalist after the Civil War, and many lesser
preachers saw their task as winning souls and preparing them for Christ's second coming. Beecher's
dream was fading. Millennial hopes dissolved in the hatreds of the war, the strikes and financial
panics of the 1870s and 1880s and the rise of an urban world that rejected Christian values, and
Evangelicals, especially in the urban North, turned more and more to speculation about the future
and to the inner life. Two statements measure the change. In 1873 Theodore Woolsey, retired president
of Yale, told the Evangelical Alliance that America is a Christian country in that the vast majority
believe in Christ and the gospel, Christian influences are universal and the civilization is built
on that foundation. In 1924 H.~L.~Mencken defined Christendom as that part of the world in which, if
a man stands up and solemnly swears that he is a Christian, all his hearers laugh. Beecher would
never have believed it.

\begin{keyfacts}
\keypeople{Lyman Beecher; James McGready; Charles G.~Finney; Theodore Weld; Harriet Beecher Stowe;
Richard Furman; Abraham Lincoln; Charles Darwin; Dwight L.~Moody; Theodore Woolsey; H.~L.~Mencken.}
\keydates{%
  \kd{1619}{First African slaves landed at Jamestown, Virginia}
  \kd{1800}{First camp meeting, at Gasper River, Kentucky}
  \kd{1835}{Lyman Beecher's \emph{A Plea for the West}}
  \kd{1837}{Theodore Weld's \emph{The Bible against Slavery}}
  \kd{1859}{Darwin's \emph{The Origin of Species}}
  \kd{1873}{Theodore Woolsey calls America a Christian country}
  \kd{1924}{H.~L.~Mencken quips that in Christendom a professing Christian is laughed at}}
\keyterms{Christian America; ``righteous empire''; frontier; voluntary society; revival; Second Great
Awakening; camp meeting; the ``invisible institution''; abolitionism; millennialism; natural
selection; higher criticism.}
\end{keyfacts}
